\documentclass[12pt,oneside]{book}

%============================================================
% PACKAGES
%============================================================

%--------------------
% Page + font
%--------------------
\usepackage[margin=1in]{geometry}
\usepackage[T1]{fontenc}
\usepackage[utf8]{inputenc}
\usepackage{lmodern}
\usepackage{microtype}

%--------------------
% Math
%--------------------
\usepackage{amsmath, amssymb, amsfonts, amsthm, mathtools}
\usepackage{bm}
\usepackage{bbm}
\usepackage{mathrsfs}
\usepackage{dsfont}

%--------------------
% Tables + figures
%--------------------
\usepackage{graphicx}
\usepackage{float}
\usepackage{booktabs}
\usepackage{array}
\usepackage{xltabular}
\usepackage{multirow}
\usepackage{caption}
\usepackage{subcaption}

%--------------------
% Lists
%--------------------
\usepackage{enumitem}
\setlist[itemize]{topsep=2pt,itemsep=2pt}
\setlist[enumerate]{topsep=2pt,itemsep=2pt}

%--------------------
% Colors + hyperlinks
%--------------------
\usepackage[dvipsnames]{xcolor}
\usepackage[
    colorlinks=true,
    linkcolor=MidnightBlue,
    citecolor=MidnightBlue,
    urlcolor=MidnightBlue
]{hyperref}
\usepackage[nameinlink,capitalise]{cleveref}

%--------------------
% Algorithms
%--------------------
\usepackage{algorithm}
\usepackage{algpseudocode}

%--------------------
% Code blocks
%--------------------
\usepackage{listings}

\lstdefinestyle{codestyle}{
    basicstyle=\ttfamily\small,
    breaklines=true,
    frame=single,
    numbers=left,
    numberstyle=\tiny,
    keywordstyle=\color{MidnightBlue},
    commentstyle=\color{ForestGreen},
    stringstyle=\color{BrickRed},
    showstringspaces=false,
    tabsize=4
}

\lstset{style=codestyle}

%--------------------
% Section styling
%--------------------
\usepackage{titlesec}
\usepackage{tocloft}

\setcounter{tocdepth}{2}
\setcounter{secnumdepth}{3}

%============================================================
% THEOREMS
%============================================================

\numberwithin{equation}{section}

\theoremstyle{definition}
\newtheorem{definition}{Definition}[section]
\newtheorem{example}{Example}[section]
\newtheorem{exercise}{Exercise}[section]
\newtheorem{remark}{Remark}[section]

\theoremstyle{plain}
\newtheorem{theorem}{Theorem}[section]
\newtheorem{lemma}{Lemma}[section]
\newtheorem{proposition}{Proposition}[section]
\newtheorem{corollary}{Corollary}[section]

\theoremstyle{remark}
\newtheorem*{note}{Note}

%============================================================
% CUSTOM COMMANDS
%============================================================

% Sets
\newcommand{\R}{\mathbb{R}}
\newcommand{\N}{\mathbb{N}}
\newcommand{\Z}{\mathbb{Z}}
\newcommand{\Q}{\mathbb{Q}}
\newcommand{\C}{\mathbb{C}}

% Probability
\newcommand{\E}{\mathbb{E}}
\newcommand{\Var}{\mathrm{Var}}
\newcommand{\Cov}{\mathrm{Cov}}
\newcommand{\Corr}{\mathrm{Corr}}
\newcommand{\Prob}{\mathbb{P}}

\newcommand{\ind}{\mathbbm{1}}
\newcommand{\iid}{\overset{\text{iid}}{\sim}}
\newcommand{\convd}{\overset{d}{\to}}
\newcommand{\convp}{\overset{p}{\to}}
\newcommand{\as}{\overset{a.s.}{\to}}

% Operators
\newcommand{\given}{\,\middle|\,}
\newcommand{\argmin}{\operatorname*{arg\,min}}
\newcommand{\argmax}{\operatorname*{arg\,max}}

% Linear algebra / analysis
\newcommand{\norm}[1]{\left\lVert #1 \right\rVert}
\newcommand{\abs}[1]{\left\lvert #1 \right\rvert}
\newcommand{\inner}[2]{\left\langle #1, #2 \right\rangle}

\newcommand{\dd}{\,\mathrm{d}}
\newcommand{\grad}{\nabla}
\newcommand{\Hess}{\nabla^2}

\newcommand{\tr}{\operatorname{tr}}
\newcommand{\rank}{\operatorname{rank}}
\newcommand{\diag}{\operatorname{diag}}

% Distributions
\newcommand{\Normal}{\mathcal{N}}
\newcommand{\Bern}{\mathrm{Bernoulli}}
\newcommand{\Bin}{\mathrm{Binomial}}
\newcommand{\Pois}{\mathrm{Poisson}}
\newcommand{\Gam}{\mathrm{Gamma}}
\newcommand{\BetaDist}{\mathrm{Beta}}
\newcommand{\Exp}{\mathrm{Exponential}}
\newcommand{\Dir}{\mathrm{Dirichlet}}
\newcommand{\InvGam}{\mathrm{Inv\text{-}Gamma}}

%============================================================
% TITLE INFO
%============================================================
\title{EN.553.724 Probabilistic Machine Learning}
\author{Jonathan Y. Ma}
\date{June 2026}

%============================================================
% DOCUMENT
%============================================================

\begin{document}

\maketitle
\tableofcontents

\newpage
\chapter{Introduction}

\section{Probabilistic Modeling and Learning}

\subsection{Classical vs.\ Modern Methods}

Classical probabilistic models (e.g.\ graphical models) and modern deep generative models differ primarily in how structure is specified.

\begin{itemize}
\item \textbf{Classical approaches}
\begin{itemize}
\item Leverage domain knowledge to specify model structure.
\item Work well when the model is correctly specified.
\item Provide interpretable probabilistic quantities and uncertainty quantification.
\end{itemize}

\item \textbf{Modern approaches}
\begin{itemize}
\item Learn highly complex distributions directly from data.
\item Require large datasets.
\item Often lack interpretability.
\end{itemize}

\item \textbf{Goal}
\[
\text{combine structured probabilistic modeling with flexible function approximation.}
\]
\end{itemize}


\subsection{Learning a Function}

Model:
\[
y = f(x) + \varepsilon, \quad \E[\varepsilon | x] = 0.
\]

Given data $\{(x_i,y_i)\}_{i=1}^N$, choose $f_\theta$.

\[
\hat{\theta} = \argmin_\theta \frac{1}{N}\sum_{i=1}^N L(f_\theta(x_i),y_i).
\]


\subsection{Maximum Likelihood Interpretation}

Assume:
\[
\varepsilon \sim \Normal(0,I).
\]

Then:
\[
p(y_i | x_i,\theta)
\propto \exp\left(-\frac{1}{2}\norm{y_i - f_\theta(x_i)}^2\right).
\]

\begin{proposition}
Minimizing squared loss is equivalent to maximum likelihood estimation.
\[
\argmax_\theta p(y_{1:N} | x_{1:N})
=
\argmin_\theta \sum_{i=1}^N \norm{y_i - f_\theta(x_i)}^2.
\]
\end{proposition}


\subsection{Learning a Distribution}

Given $x_1,\dots,x_N \iid \Prob$:

\[
\hat{\theta} = \argmax_\theta \sum_{i=1}^N \log p_\theta(x_i)
=
\argmin_\theta \frac{1}{N}\sum_{i=1}^N -\log p_\theta(x_i).
\]


\section{Measure-Theoretic Foundations}

\subsection{Discrete and Continuous Probability}

\begin{itemize}
\item Discrete:
\[
\Prob(A) = \sum_{x \in A} p(x)
\]

\item Continuous:
\[
\Prob(A) = \int_A p(x)\,\dd x
\]
\end{itemize}

Unified via:
\[
\Prob(A) = \int_A p(x)\,\lambda(\dd x).
\]


\subsection{Probability Spaces}

\begin{definition}
A measurable space is $(\Omega,\mathcal{F})$.
\end{definition}

\begin{definition}
A probability measure $\Prob$ satisfies:
\begin{itemize}
\item $\Prob(A) \ge 0$
\item $\Prob(\Omega) = 1$
\item countable additivity
\end{itemize}
\end{definition}

If a density exists:
\[
\Prob(A) = \int_A p(x)\,\lambda(\dd x).
\]


\subsection{Expectation}

\[
\E[f] = \int f(x)\,\Prob(\dd x)
\]

If density exists:
\[
\E[f] = \int f(x)p(x)\,\lambda(\dd x).
\]


\subsection{Radon--Nikodym Derivative}

\begin{definition}
$\Prob \ll Q$ if $Q(A)=0 \Rightarrow \Prob(A)=0$.
\end{definition}

\begin{theorem}[Radon--Nikodym]
If $\Prob \ll Q$, then
\[
\int f \dd \Prob = \int f \frac{\dd \Prob}{\dd Q} \dd Q.
\]
\end{theorem}

If densities exist:
\[
\frac{\dd \Prob}{\dd Q}(x) = \frac{p(x)}{q(x)}.
\]


\subsection{Random Variables}

\begin{definition}
A random variable is a measurable function $X:\Omega \to \R$.
\end{definition}

Distribution:
\[
\Prob_X(A) = \Prob(X \in A).
\]


\subsection{Conditional Probability}

\[
\Prob(A | B) = \frac{\Prob(A \cap B)}{\Prob(B)}
\]

\[
p(x|y) = \frac{p(x,y)}{p(y)}
\]

\[
p(x|y) = \frac{p(x)p(y|x)}{p(y)}
\]


\subsection{Beta-Binomial Model}

\[
\theta \sim \BetaDist(\alpha,\beta), \quad X|\theta \sim \Bin(m+n,\theta).
\]

Posterior:
\[
\theta | X \sim \BetaDist(\alpha+m,\beta+n).
\]


\subsection{Conditional Independence}

\begin{definition}
$X \perp Y | Z$ if
\[
p(x,y|z) = p(x|z)p(y|z).
\]
\end{definition}


\section{Divergences and Information}

\subsection{Total Variation}

\begin{definition}
\[
\mathrm{TV}(\Prob,Q) = \frac{1}{2}\int |p-q|\,\lambda(\dd x).
\]
\end{definition}

\begin{proposition}
\[
\mathrm{TV}(\Prob,Q) = \max_{A} |\Prob(A)-Q(A)|.
\]
\end{proposition}


\subsection{$f$-Divergence}

\begin{definition}
\[
D_f(\Prob \| Q)
=
\E_Q\left[f\left(\frac{p}{q}\right)\right].
\]
\end{definition}

\[
\mathrm{KL}(\Prob \| Q)
=
\int p\log\frac{p}{q}.
\]


\subsection{Non-Negativity}

\begin{theorem}[Jensen]
\[
\E[f(X)] \ge f(\E[X]).
\]
\end{theorem}

\begin{proposition}
\[
D_f(\Prob \| Q) \ge 0.
\]
\end{proposition}

\begin{proof}
\[
\E_Q f\left(\frac{p}{q}\right)
\ge
f\left(\E_Q \frac{p}{q}\right)
= f(1)=0.
\]
\end{proof}


\subsection{Entropy and KL}

\begin{definition}
\[
H(\Prob) = -\int p\log p.
\]
\end{definition}

\[
\mathrm{KL}(\Prob \| Q)
=
-H(\Prob) - \E_\Prob[\log q].
\]


\subsection{MLE and KL}

\begin{proposition}
\[
\E_\Prob[-\log p_\theta(x)]
=
\mathrm{KL}(\Prob \| \Prob_\theta) + H(\Prob).
\]
\end{proposition}

\[
\Rightarrow
\text{MLE} = \argmin_\theta \mathrm{KL}(\Prob \| \Prob_\theta).
\]


\subsection{KL Direction}

\[
\mathrm{KL}(\Prob \| Q) \quad \text{mass-covering}
\]

\[
\mathrm{KL}(Q \| \Prob) \quad \text{mode-seeking}
\]


\subsection{Invariance}

\begin{proposition}
For invertible $g$:
\[
D_f(\Prob \| Q)
=
D_f(g_\#\Prob \| g_\# Q).
\]
\end{proposition}

\section{Optimal Transport and Integral Probability Metrics}

\subsection{Couplings and Wasserstein Distance}

\begin{definition}
Let $(\Omega,\mathcal{F})$ be a measurable space with metric $d$, and let $\mu,\nu$ be probability measures on $\Omega$.

A \emph{coupling} of $\mu$ and $\nu$ is a probability measure $\gamma$ on $\Omega \times \Omega$ such that:
\[
\gamma(A \times \Omega) = \mu(A), \quad
\gamma(\Omega \times A) = \nu(A).
\]
\end{definition}

Interpretation: a coupling specifies a joint distribution $(X,Y)$ with marginals $X \sim \mu$, $Y \sim \nu$.

\begin{definition}
The Wasserstein-$p$ distance is defined as
\[
W_p(\mu,\nu)
=
\left(
\inf_{\gamma \in \Gamma(\mu,\nu)}
\E_{(X,Y)\sim \gamma}[d(X,Y)^p]
\right)^{1/p},
\]
where $\Gamma(\mu,\nu)$ is the set of all couplings of $\mu,\nu$.
\end{definition}

\begin{remark}
$W_1$ is often called the \emph{earth mover's distance}, since it measures the minimal cost of transporting mass from $\mu$ to $\nu$.
\end{remark}


\subsection{Integral Probability Metrics}

\begin{definition}
Let $\mathcal{F}$ be a class of functions $f : \Omega \to \R$. The integral probability metric (IPM) is
\[
d_{\mathcal{F}}(\nu,\mu)
=
\sup_{f \in \mathcal{F}}
\left|
\E_\nu[f] - \E_\mu[f]
\right|.
\]
\end{definition}

\begin{lemma}
Let $\mathcal{F} = \{f : \norm{f}_{\infty} \le 1\}$. Then
\[
\mathrm{TV}(\mu,\nu) = \frac{1}{2} d_{\mathcal{F}}(\nu,\mu).
\]
\end{lemma}

\begin{proof}
For any $f$ with $\norm{f}_{\infty} \le 1$,
\[
|\E_\nu[f] - \E_\mu[f]|
=
\left| \int f(x)(\nu - \mu)(\dd x) \right|
\le \int |f(x)|\,|\nu - \mu|(\dd x)
\le \int |\nu - \mu|(\dd x).
\]

Taking the supremum gives
\[
d_{\mathcal{F}}(\nu,\mu) \le 2\,\mathrm{TV}(\mu,\nu).
\]

Equality is achieved by choosing
\[
f(x) = \mathrm{sign}(p(x) - q(x)),
\]
which yields
\[
\E_\nu[f] - \E_\mu[f] = \int |p(x)-q(x)| \dd x.
\]
\end{proof}


\subsection{Kantorovich--Rubinstein Duality}

\begin{theorem}[Kantorovich--Rubinstein]
Let $\mathcal{F}$ be the set of 1-Lipschitz functions on $\R^d$. Then
\[
W_1(\mu,\nu) = d_{\mathcal{F}}(\nu,\mu).
\]
\end{theorem}

\begin{proof}[Sketch]
For any coupling $\gamma$ and any 1-Lipschitz $f$,
\[
|f(x) - f(y)| \le d(x,y).
\]

Thus
\[
\E_\nu[f] - \E_\mu[f]
=
\E_\gamma[f(Y) - f(X)]
\le \E_\gamma[d(X,Y)].
\]

Taking infimum over $\gamma$ gives
\[
d_{\mathcal{F}}(\nu,\mu) \le W_1(\mu,\nu).
\]

The reverse inequality follows by constructing an optimal dual potential.
\end{proof}


\subsection*{Remarks on Metrics vs Divergences}

\begin{itemize}
\item Total variation and Wasserstein are \emph{metrics}.
\item $f$-divergences (e.g.\ KL) are not metrics:
\begin{itemize}
\item not symmetric
\item do not satisfy triangle inequality
\end{itemize}
\item IPMs provide a unifying view of distances via function classes.
\end{itemize}

\section{Maximum Entropy and Exponential Families}

\subsection{Maximum Entropy Principle}

Let $P$ be an unknown distribution on $\Omega$. Suppose we are given moment constraints:
\[
\E_P[\phi_i(X)] = \mu_i, \quad i = 1,\dots,d,
\]
or equivalently
\[
\E_P[\phi(X)] = \mu \in \R^d,
\]
where $\phi : \Omega \to \R^d$.

\begin{remark}
These constraints alone do not uniquely determine $P$.
\end{remark}

A natural principle is to choose the \emph{least informative} distribution satisfying the constraints.

\begin{definition}
The maximum entropy distribution (with respect to a reference measure $\nu$) is
\[
\argmax_{P}
\; H_\nu(P)
\quad \text{s.t.} \quad \E_P[\phi] = \mu,
\]
where
\[
H_\nu(P) = -\int p(x)\log p(x)\,\nu(\dd x).
\]
\end{definition}


\subsection*{Derivation via Lagrange Multipliers}

Assume $\Omega$ is finite for simplicity, and write $p(x) = P(X=x)$.

We maximize
\[
-\sum_{x \in \Omega} p(x)\log p(x)
\]
subject to:
\[
\sum_x p(x) = 1, \quad
\sum_x p(x)\phi_i(x) = \mu_i.
\]

Define the Lagrangian:
\[
\mathcal{L}(p,\lambda_0,\lambda)
=
-\sum_x p(x)\log p(x)
+ \lambda_0\left(\sum_x p(x)-1\right)
+ \sum_{i=1}^d \lambda_i\left(\sum_x p(x)\phi_i(x)-\mu_i\right).
\]

Taking derivatives:
\[
\frac{\partial \mathcal{L}}{\partial p(x)}
=
-\log p(x) - 1 + \lambda_0 + \sum_{i=1}^d \lambda_i \phi_i(x).
\]

Setting to zero:
\[
\log p(x)
=
\lambda_0 - 1 + \sum_{i=1}^d \lambda_i \phi_i(x).
\]

Thus
\[
p(x)
\propto
\exp\left(\sum_{i=1}^d \lambda_i \phi_i(x)\right)
=
\exp(\langle \theta, \phi(x)\rangle),
\]
where $\theta = (\lambda_1,\dots,\lambda_d)$.

This gives the exponential family form.


\subsection{Exponential Family}

\begin{definition}
Let $(\mathcal{X},\nu)$ be a measure space and $\phi : \mathcal{X} \to \R^d$.

The exponential family is
\[
p_\theta(x)
=
\frac{\exp(\langle \theta,\phi(x)\rangle)}{Z(\theta)}
=
\exp\big(\langle \theta,\phi(x)\rangle - A(\theta)\big),
\]
where
\[
Z(\theta) = \int \exp(\langle \theta,\phi(x)\rangle)\,\nu(\dd x),
\quad
A(\theta) = \log Z(\theta).
\]
\end{definition}

\begin{remark}
\begin{itemize}
\item $\theta$ is the natural parameter
\item $\phi(x)$ are sufficient statistics
\item $A(\theta)$ is the log-partition function
\end{itemize}
\end{remark}


\subsubsection{Example: Gaussian from Exponential Family}

Let $\phi(x) = (x, x^2)$ on $\R$.

Then
\[
p_\theta(x)
\propto
\exp(\theta_1 x + \theta_2 x^2).
\]

Completing the square:
\[
\theta_1 x + \theta_2 x^2
=
\theta_2 \left(x + \frac{\theta_1}{2\theta_2}\right)^2
- \frac{\theta_1^2}{4\theta_2}.
\]

For $\theta_2 < 0$, this is a Gaussian:
\[
p_\theta(x) = \Normal(\mu,\sigma^2),
\quad
\mu = -\frac{\theta_1}{2\theta_2},
\quad
\sigma^2 = -\frac{1}{2\theta_2}.
\]


\subsection{Maximum Entropy Characterization}

\begin{theorem}
If there exists $\theta$ such that
\[
\E_{p_\theta}[\phi(X)] = \mu,
\]
then $p_\theta$ is the maximum entropy distribution among all distributions satisfying the moment constraints.
\end{theorem}

\begin{proof}
Let $P$ satisfy $\E_P[\phi]=\mu$. Then
\[
\mathrm{KL}(P \| p_\theta)
=
\E_P[\log p(x) - \log p_\theta(x)] \ge 0.
\]

Rearranging:
\[
-H(P) \ge \E_P[\log p_\theta(x)].
\]

Using exponential family form:
\[
\E_P[\log p_\theta(x)]
=
\E_P[\langle \theta,\phi(x)\rangle - A(\theta)]
=
\langle \theta, \mu\rangle - A(\theta).
\]

Thus
\[
H(P)
\le
A(\theta) - \langle \theta,\mu\rangle
=
H(p_\theta),
\]
so $p_\theta$ maximizes entropy.
\end{proof}


\subsection{Mean Parameter Mapping}

Define the mapping
\[
\nabla A(\theta) = \E_{p_\theta}[\phi(X)].
\]

\begin{definition}
Let
\[
\mathcal{M}
=
\{\mu \in \R^d : \exists P \text{ such that } \E_P[\phi(X)] = \mu\}
\]
be the set of achievable mean parameters.
\end{definition}

\begin{definition}
The statistics $\{\phi_i\}_{i=1}^d$ are \emph{minimal} if
\[
\{1,\phi_1,\dots,\phi_d\}
\text{ are linearly independent}.
\]
\end{definition}

\begin{proposition}
If $\phi$ is minimal, then:
\begin{itemize}
\item $\nabla A : \Theta \to \mathcal{M}$ is injective
\item $\nabla A : \Theta \to \mathcal{M}^\circ$ is surjective (onto the interior)
\end{itemize}
\end{proposition}

\begin{proof}[Sketch]
If $\phi$ is minimal, then $A(\theta)$ is strictly convex:
\[
\nabla^2 A(\theta) = \mathrm{Cov}_{p_\theta}(\phi(X)) \succ 0.
\]

Strict convexity implies injectivity of $\nabla A$.

Surjectivity onto the interior follows from convex duality and properties of cumulant generating functions.
\end{proof}


\subsection{Convexity and Geometry}

\begin{proposition}
The log-partition function satisfies:
\[
\nabla^2 A(\theta) = \mathrm{Cov}_{p_\theta}(\phi(X)) \succeq 0.
\]
\end{proposition}

\begin{remark}
\begin{itemize}
\item $A(\theta)$ is convex
\item strictly convex if $\phi$ is minimal
\item exponential families form a convex dual pair:
\[
\theta \leftrightarrow \mu = \nabla A(\theta)
\]
\end{itemize}
\end{remark}


\subsection*{Non-Minimal Case}

If $\phi$ is not minimal, then there exists $v \neq 0$ such that
\[
\langle v, \phi(x)\rangle = c \quad \forall x.
\]

Then
\[
p_{\theta + t v}(x)
=
p_\theta(x),
\]
so the parametrization is not identifiable, and $\nabla A$ is not injective.

\chapter{Graphical Models}

\section{Directed Graphical Models (Bayesian Networks)}

\subsection{Factorization in DAGs}

\begin{definition}
Let $G = (V,E)$ be a directed acyclic graph (DAG), and let $\{X_v\}_{v \in V}$ be random variables.

A distribution $P$ \emph{factorizes according to} $G$ if
\[
p(x_V) = \prod_{v \in V} p\big(x_v \mid x_{\mathrm{pa}(v)}\big),
\]
where $\mathrm{pa}(v)$ denotes the set of parents of node $v$.
\end{definition}

\begin{remark}
Each factor $p(x_v \mid x_{\mathrm{pa}(v)})$ is a \emph{local conditional model}.
\end{remark}

\begin{definition}
A \emph{topological ordering} of a DAG is an ordering $(v_1,\dots,v_n)$ such that
\[
v_i \to v_j \implies i < j.
\]
\end{definition}

\begin{remark}
This allows a generative interpretation:
\[
X_{v_1} \to X_{v_2} \to \cdots \to X_{v_n},
\]
sampling each variable conditioned on its parents.
\end{remark}


\subsection{Independence Structure}

\begin{definition}
Let $S_v = \big(\{v\} \cup \mathrm{Desc}(v) \cup \mathrm{pa}(v)\big)^c$.

A DAG encodes the independence relations
\[
X_v \perp X_{S_v} \mid X_{\mathrm{pa}(v)}.
\]
\end{definition}

\begin{proposition}
A distribution factorizes according to $G$ if and only if it satisfies the independence relations encoded by $G$.
\end{proposition}


\subsection{Factorization $\Rightarrow$ Independence (Key Idea)}

\begin{lemma}
If $P$ factorizes according to a DAG $G$, then for any upward-closed set $U$,
\[
p(x_U) = \prod_{u \in U} p(x_u \mid x_{\mathrm{pa}(u)}).
\]
\end{lemma}

\begin{proof}[Sketch]
Sum out variables not in $U$ using the factorization; terms outside $U$ cancel due to normalization.
\end{proof}

This leads to conditional independence:
\[
X_v \perp X_{S_v} \mid X_{\mathrm{pa}(v)}.
\]


\subsection{Constructing a DAG}

Given variables $X_1,\dots,X_n$:

\begin{itemize}
\item Choose an ordering $X_1,\dots,X_n$
\item For each $k$, choose a minimal set $S \subseteq \{1,\dots,k-1\}$ such that
\[
X_k \perp X_{[k-1]\setminus S} \mid X_S
\]
\item Set $\mathrm{pa}(k) = S$
\end{itemize}

\begin{remark}
The resulting graph is not unique.
\end{remark}

\begin{remark}
Causality can guide graph construction, but cannot generally be inferred from the graph alone.
\end{remark}


\subsection{Compactness of Representation}

For $n$ variables taking values in $\{1,\dots,r\}$:

\begin{itemize}
\item Full joint distribution requires:
\[
r^n - 1 \quad \text{parameters}
\]
\item If max indegree is $d$:
\[
O(n r^{d+1})
\]
parameters suffice
\end{itemize}

\begin{remark}
Graphical models exploit conditional independence for exponential savings.
\end{remark}


\subsection{Examples}

\subsubsection{Markov Chain}

\[
X_{k+1} \perp X_1,\dots,X_{k-1} \mid X_k
\]

Factorization:
\[
p(x_1,\dots,x_n)
=
p(x_1)\prod_{k=2}^n p(x_k \mid x_{k-1})
\]

\subsubsection{$d$-th Order Markov Model}

\[
X_{k+d} \perp X_1,\dots,X_{k-1} \mid X_k,\dots,X_{k+d-1}
\]

\subsubsection{Naive Bayes}

Latent class $C$:
\[
X_k \perp X_{-k} \mid C
\]

Factorization:
\[
p(c,x_1,\dots,x_n)
=
p(c)\prod_{k=1}^n p(x_k \mid c)
\]


\subsection{Latent Variable Models}

\begin{definition}
A \emph{latent variable model} introduces hidden variables $Z$ such that
\[
p(x) = \int p(x,z)\,\dd z.
\]
\end{definition}

\begin{remark}
Latent variables simplify structure by introducing conditional independence.
\end{remark}

\subsubsection{Mixture Model}

\[
z_i \sim \mathrm{Multinomial}(\pi), \quad
x_i \mid z_i \sim \Normal(\mu_{z_i}, \Sigma_{z_i})
\]

\subsubsection{Stochastic Block Model}

\[
z_i \sim \mathrm{Multinomial}(\pi), \quad
A_{ij} \mid z_i,z_j \sim \mathrm{Bernoulli}(p_{z_i z_j})
\]

\subsubsection{Noisy-OR Model}

\[
x_j \mid z \sim \mathrm{Bernoulli}\left(1 - \exp\left(-\sum_i w_{ij} z_i\right)\right)
\]

\subsubsection{Hidden Markov Model (HMM)}

\[
z_t = A z_{t-1} + \epsilon_t, \quad
x_t = C z_t + \eta_t
\]

\begin{remark}
This is a linear Gaussian state-space model.
\end{remark}


\subsection{d-Separation}

\begin{definition}
A path in a DAG is \emph{blocked} by a set $C$ if one of the following holds:
\begin{enumerate}
\item Chain: $X \to Z \to Y$ or $X \leftarrow Z \leftarrow Y$, and $Z \in C$
\item Fork: $X \leftarrow Z \to Y$, and $Z \in C$
\item Collider: $X \to Z \leftarrow Y$, and $Z \notin C$ and no descendant of $Z$ is in $C$
\end{enumerate}
\end{definition}

\begin{definition}
$A$ and $B$ are \emph{d-separated} by $C$ if all paths between them are blocked.
\end{definition}

\begin{theorem}
If $P$ factorizes according to $G$ and $A,B$ are d-separated by $C$, then
\[
X_A \perp X_B \mid X_C.
\]
\end{theorem}


\section{Undirected Graphical Models (Markov Random Fields)}

\subsection{Factorization}

\begin{definition}
Let $G=(V,E)$ be an undirected graph. Let $\mathcal{C}$ be the set of maximal cliques.

A distribution factorizes over $G$ if
\[
p(x) \propto \exp(-E(x)),
\quad
E(x) = \sum_{C \in \mathcal{C}} \phi_C(x_C),
\]
or equivalently
\[
p(x) = \prod_{C \in \mathcal{C}} \psi_C(x_C).
\]
\end{definition}

\begin{remark}
$\psi_C = e^{-\phi_C}$ are called \emph{clique potentials}.
\end{remark}


\subsection{Markov Properties}

\begin{definition}
$P$ satisfies:
\begin{itemize}
\item Global Markov property:
\[
X_A \perp X_B \mid X_C
\quad \text{if } C \text{ separates } A,B
\]
\item Local Markov property:
\[
X_v \perp X_{V \setminus (\{v\}\cup N(v))} \mid X_{N(v)}
\]
\item Pairwise Markov property:
\[
X_v \perp X_w \mid X_{V \setminus \{v,w\}}, \quad (v,w)\notin E
\]
\end{itemize}
\end{definition}

\begin{remark}
Under positivity conditions:
\[
\text{global} \iff \text{local} \iff \text{pairwise}
\]
\end{remark}


\subsection{Energy-Based Models}

\begin{definition}
An energy-based model defines
\[
p_\theta(x) \propto e^{-E_\theta(x)}.
\]
\end{definition}

\begin{remark}
\begin{itemize}
\item Lower energy $\Rightarrow$ higher probability
\item $E_\theta(x) = \langle \theta, \phi(x)\rangle$ recovers exponential families
\end{itemize}
\end{remark}


\subsection{Example: Ising Model}

\[
p(x)
\propto
\exp\left(
\sum_i h_i x_i + \sum_{(i,j)\in E} J_{ij} x_i x_j
\right),
\quad x_i \in \{-1,1\}
\]

Conditional distribution:
\[
p(x_i = 1 \mid x_{-i})
=
\frac{\exp(h_i + \sum_{j\in N(i)} J_{ij} x_j)}
{\exp(h_i + \sum_{j\in N(i)} J_{ij} x_j) + \exp(-h_i - \sum_{j\in N(i)} J_{ij} x_j)}.
\]

\begin{remark}
Unlike DAGs, sampling is generally difficult due to the global partition function.
\end{remark}

\section{Additional Properties of Independence}

\subsection*{Intersection Property}

\begin{proposition}
Suppose $p(x,y,z,w)$ has positive density with respect to a product measure. Then the \emph{intersection property} holds:
\[
\begin{cases}
X \perp Y \mid Z, W, \\
X \perp W \mid Z, Y
\end{cases}
\quad \Rightarrow \quad
X \perp (Y,W) \mid Z.
\]
\end{proposition}

\begin{remark}
This property does \emph{not} hold in general without positivity.
\end{remark}

\begin{proposition}
If the intersection property holds, then the global, local, and pairwise Markov properties are equivalent.
\end{proposition}


\subsection*{Hammersley--Clifford Theorem}

\begin{theorem}[Hammersley--Clifford]
Let $G=(V,E)$ be an undirected graph. Suppose $P$ has positive density.

Then:
\[
P \text{ satisfies the Markov properties}
\;\Longleftrightarrow\;
P \text{ factorizes over } G.
\]
\end{theorem}

\begin{remark}
Positivity is essential: without it, independence does not imply factorization.
\end{remark}


\subsection*{Canonical Representation}

\begin{theorem}
Any strictly positive distribution can be written as
\[
p(x) = \exp\left( \sum_{S \subseteq V} \xi_S(x_S) \right),
\]
where
\[
\xi_S(x_S) = \sum_{Z \subseteq S} (-1)^{|S|-|Z|} \log p(x_Z, x_{S^c}=0).
\]
\end{theorem}

\begin{remark}
\begin{itemize}
\item This is an inclusion--exclusion type expansion
\item If $P$ factorizes over $G$, then $\xi_S = 0$ for non-cliques $S$
\end{itemize}
\end{remark}


\section{Factor Graphs}

\begin{definition}
A \emph{factor graph} is a bipartite graph with:
\begin{itemize}
\item variable nodes $V$
\item factor nodes $F$
\end{itemize}
and edges only between variables and factors.

A distribution factorizes as
\[
p(x_V) = \prod_{a \in F} \varphi_a(x_{N(a)}),
\]
where $N(a)$ are variables connected to factor $a$.
\end{definition}

\begin{remark}
Factor graphs unify:
\begin{itemize}
\item DAG factorizations
\item MRF factorizations
\end{itemize}
\end{remark}


\subsection{Example: Restricted Boltzmann Machine (RBM)}

Let visible units $v \in \{0,1\}^{d_v}$ and hidden units $h \in \{0,1\}^{d_h}$.

\[
p(v,h) \propto \exp\left( \langle v,Wh\rangle + \langle b,v\rangle + \langle a,h\rangle \right)
\]

\begin{remark}
Key property:
\[
p(h \mid v) = \prod_j p(h_j \mid v),
\quad
p(v \mid h) = \prod_i p(v_i \mid h)
\]
\end{remark}

\begin{remark}
This conditional independence enables efficient Gibbs sampling.
\end{remark}


\subsection{Example: Coding Theory (LDPC)}

\[
p(x,z)
\propto
\prod_i \mathbbm{1}\{H_i z = 0\}
\prod_j p(x_j \mid z_j)
\]

\begin{remark}
Inference corresponds to belief propagation.
\end{remark}


\section{Conditional Random Fields}

\begin{definition}
Let $(V \cup W, E)$ be an undirected graph.

A \emph{conditional random field (CRF)} is a conditional distribution
\[
P(X_W \mid X_V)
\propto
\prod_{C \in \mathcal{C}} \psi_C(x_C),
\]
where cliques $C$ involve variables in $W$.
\end{definition}

\begin{remark}
CRFs model conditional structure without specifying a joint model over $(V,W)$.
\end{remark}


\section{Directed $\to$ Undirected Conversion}

\begin{definition}
The \emph{moralized graph} $\mathcal{M}[G]$ of a DAG $G$ is obtained by:
\begin{itemize}
\item adding edges between all pairs of parents of each node
\item dropping edge directions
\end{itemize}
\end{definition}

\begin{proposition}
If $P$ factorizes according to a DAG $G$, then it factorizes according to the undirected graph $\mathcal{M}[G]$.
\end{proposition}


\section{Independence via Moralization}

\begin{theorem}
Let $G$ be a DAG and $A,B,C$ disjoint sets. The following are equivalent:

\begin{enumerate}
\item $A$ and $B$ are d-separated by $C$ in $G$
\item $A$ and $B$ are separated by $C$ in the moralized ancestral graph
\end{enumerate}

Moreover, if $P$ factorizes according to $G$, then
\[
X_A \perp X_B \mid X_C.
\]
\end{theorem}

\begin{definition}
The \emph{ancestral graph} is the subgraph induced by $A,B,C$ and their ancestors.
\end{definition}


\section{Markov Blanket}

\begin{definition}
A set $U$ is a \emph{Markov blanket} for $v$ if
\[
X_v \perp X_{V \setminus (U \cup \{v\})} \mid X_U.
\]
\end{definition}

\begin{proposition}
\begin{enumerate}
\item In an undirected graph, $N(v)$ is a Markov blanket.
\item In a DAG, the Markov blanket is:
\[
\mathrm{pa}(v) \cup \mathrm{ch}(v) \cup \mathrm{co\text{-}pa}(v).
\]
\end{enumerate}
\end{proposition}

\begin{remark}
The Markov blanket gives the minimal conditioning set needed to isolate a node.
\end{remark}


\section{Learning Graph Structure (Undirected Case)}

\begin{proposition}
Assume $P$ is positive. Define graph $G$ by:
\[
(v,w) \in E
\;\Longleftrightarrow\;
X_v \not\perp X_w \mid X_{V \setminus \{v,w\}}.
\]

Then $G$ is the unique minimal graph for which $P$ satisfies the pairwise Markov property.
\end{proposition}


\section*{Comparison: Directed vs Undirected Models}

\begin{itemize}
\item Directed models:
\begin{itemize}
\item capture causal / generative structure
\item sampling is straightforward
\item require ordering / domain knowledge
\end{itemize}

\item Undirected models:
\begin{itemize}
\item symmetric relationships
\item easier structure learning
\item inference often harder (partition function)
\end{itemize}
\end{itemize}

\section{Inference Problems}

\subsection{Core Tasks}

Given a probabilistic model $p_\theta(x)$:

\begin{itemize}
\item \textbf{Sampling:}
\[
x \sim p_\theta(x)
\]

\item \textbf{Marginals:}
\[
p_\theta(x_i) = \int p_\theta(x)\,\nu^{\otimes (n-1)}(dx_{-i})
\]

\item \textbf{MAP (maximum a posteriori):}
\[
x^* = \argmax_x p_\theta(x)
\]

\item \textbf{Partition function:}
\[
Z(\theta) = \int \psi_\theta(x)\,\nu^{\otimes n}(dx)
\]
\end{itemize}

\begin{remark}
These tasks are tightly coupled: computing one often requires solving the others.
\end{remark}


\subsection*{Latent Variable Models}

For joint models:
\[
p_\theta(x,z) \propto \psi_\theta(x,z)
\]

posterior inference requires:
\[
p_\theta(z|x) = \frac{\psi_\theta(x,z)}{Z(\theta,x)}, \quad
Z(\theta,x) = \int \psi_\theta(x,z)\,\nu_Z(dz)
\]

\begin{itemize}
\item Sampling: $z \sim p_\theta(z|x)$
\item Marginals: $p_\theta(z_i|x)$
\item MAP: $\argmax_z p_\theta(z|x)$
\end{itemize}


\section{Computational Complexity}

\subsection*{Complexity Classes}

\begin{itemize}
\item $O(n)$, $O(n\log n)$, $O(n^c)$: polynomial time
\item $O(2^n)$: exponential time
\end{itemize}

\begin{definition}
A \emph{decision problem} outputs YES/NO.

\begin{itemize}
\item \textbf{P}: solvable in polynomial time
\item \textbf{NP}: solutions verifiable in polynomial time
\item \textbf{NP-complete}: hardest problems in NP
\item \textbf{NP-hard}: at least as hard as NP-complete
\end{itemize}
\end{definition}

\begin{example}
3-SAT: given a Boolean formula, does a satisfying assignment exist?
\end{example}


\subsection*{Counting Problems}

\begin{definition}
$\#\text{P}$: counting versions of NP problems (e.g., number of satisfying assignments).
\end{definition}

\begin{remark}
Computing partition functions and marginals is typically $\#\text{P}$-hard.
\end{remark}


\begin{theorem}
Computing partition functions or marginals in general graphical models is $\#\text{P}$-hard.
\end{theorem}

\begin{remark}
Even approximate counting (e.g., $\#\text{SAT}$) can be NP-hard.
\end{remark}


\section{Approximation Algorithms}

\begin{definition}
A \emph{polynomial-time randomized approximation scheme (PRAS)} for $Z$ outputs $\widehat{Z}$ such that
\[
e^{-\epsilon} Z \le \widehat{Z} \le e^{\epsilon} Z
\]
in polynomial time, with high probability.
\end{definition}

\begin{definition}
\begin{itemize}
\item \textbf{FPRAS}: runtime polynomial in $n$ and $1/\epsilon$
\item \textbf{PTAS}: deterministic approximation scheme
\end{itemize}
\end{definition}


\section*{Equivalence: Sampling, Marginals, Partition Function}

\begin{definition}
A family of problems is \emph{self-reducible} if fixing variables reduces the problem to a smaller instance.
\end{definition}

\begin{theorem}[informal]
For distributions $p_\theta(x) \propto \psi_\theta(x)$ on finite spaces:

\begin{enumerate}
\item Approximate sampling $\Rightarrow$ approximate marginals
\item Approximate marginals $\Rightarrow$ approximate sampling (if self-reducible)
\item Marginals $\Leftrightarrow$ partition function
\end{enumerate}
\end{theorem}


\begin{theorem}
If we can sample within TV distance $\epsilon$ in $\mathrm{poly}(n,1/\epsilon)$ time, then we can approximate marginals via empirical averages:
\[
\widehat{p}(x_i = a) = \frac{1}{N} \sum_{k=1}^N \ind\{x_i^{(k)} = a\}
\]
\end{theorem}

\begin{proof}
By Hoeffding’s inequality:
\[
\Prob\left( |\widehat{p} - p| \ge \delta \right)
\le 2 \exp(-2N\delta^2)
\]
so $N = O(1/\delta^2)$ suffices.
\end{proof}


\section{Exact Inference}

\subsection{Variable Elimination}

\begin{algorithm}
\caption{Variable Elimination}
\begin{algorithmic}
\State Input: $p(x) \propto \prod_a \psi_a(x_{\partial a})$
\State Choose elimination order
\For{variable $i$}
    \State Multiply all factors involving $i$
    \State Sum out $i$
    \State Introduce new factor
\EndFor
\State Output desired marginal
\end{algorithmic}
\end{algorithm}

\begin{remark}
Intermediate factors grow in size → exponential complexity.
\end{remark}


\subsection*{Treewidth}

\begin{definition}
The \emph{treewidth} $\mathrm{tw}(G)$ is the size of the largest intermediate factor (minus one) in optimal elimination ordering.
\end{definition}

\begin{remark}
Complexity of variable elimination:
\[
O(n \cdot \text{state}^{\mathrm{tw}(G)+1})
\]
\end{remark}


\section{Belief Propagation}

\subsection{Message Passing}

Messages:

\begin{itemize}
\item Variable $\to$ factor:
\[
m_{i \to a}(x_i)
\propto \prod_{b \in \partial i \setminus a} m_{b \to i}(x_i)
\]

\item Factor $\to$ variable:
\[
m_{a \to i}(x_i)
\propto
\sum_{x_{\partial a \setminus i}}
\psi_a(x_{\partial a})
\prod_{j \in \partial a \setminus i} m_{j \to a}(x_j)
\]
\end{itemize}


\subsection*{Marginals}

\[
p(x_i)
\propto \prod_{a \in \partial i} m_{a \to i}(x_i)
\]

\[
p(x_{\partial a})
\propto
\psi_a(x_{\partial a})
\prod_{j \in \partial a} m_{j \to a}(x_j)
\]


\subsection{Properties}

\begin{itemize}
\item Exact on trees
\item Equivalent to dynamic programming
\item Linear-time in number of nodes (for trees)
\end{itemize}

\begin{remark}
On loopy graphs:
\begin{itemize}
\item becomes \emph{loopy belief propagation}
\item no guarantees, but often works well
\end{itemize}
\end{remark}


\section{Extensions}

\begin{itemize}
\item Conditioning: add indicator factors
\item Continuous variables: approximate messages (e.g., Gaussian)
\item Junction tree algorithm: exact inference via tree decomposition
\end{itemize}


\section{Takeaways}

\begin{itemize}
\item Exact inference is exponential in general
\item Tree structure enables efficient computation
\item Message passing = dynamic programming on graphs
\item Sampling, marginals, and partition function are deeply connected
\end{itemize}

\subsection*{Correctness of Belief Propagation (Trees)}

We formalize what the messages compute.

Let $T_{b \to i}$ denote the subtree rooted at $b$ when edge $(b,i)$ is removed.

\begin{proposition}
For a tree-structured factor graph:
\[
Z_{b \to i}(x_i)
=
\sum_{x_{V(T_{b \to i}) \setminus i}}
\prod_{c \in F(T_{b \to i})} \psi_c(x_{\partial c}),
\]
and similarly
\[
Z_{a \to j}(x_j)
=
\sum_{x_{V(T_{a \to j}) \setminus j}}
\prod_{c \in F(T_{a \to j})} \psi_c(x_{\partial c}).
\]
\end{proposition}

\begin{proof}
By induction on the size of the subtree.

\textbf{Base case:} leaf node. Only one factor contributes, and the expression holds trivially.

\textbf{Inductive step:} assume the claim holds for all subtrees of depth $\le k$.  
Then for a larger subtree,
\[
Z_{a \to i}(x_i)
=
\sum_{x_{\partial a \setminus i}}
\psi_a(x_{\partial a})
\prod_{j \in \partial a \setminus i} Z_{j \to a}(x_j),
\]
which matches the factor-to-variable message update.

Thus the recursion exactly computes subtree marginals.
\end{proof}


\subsection*{Max-Product Belief Propagation (MAP)}

Replace sums by maxima.

\begin{itemize}
\item Variable $\to$ factor:
\[
M_{i \to a}(x_i)
=
\prod_{b \in \partial i \setminus a} M_{b \to i}(x_i)
\]

\item Factor $\to$ variable:
\[
M_{a \to i}(x_i)
=
\max_{x_{\partial a \setminus i}}
\psi_a(x_{\partial a})
\prod_{j \in \partial a \setminus i} M_{j \to a}(x_j)
\]
\end{itemize}

\begin{remark}
This computes max-marginals:
\[
M_i(x_i) \propto \max_{x_{-i}} p(x)
\]
\end{remark}


\begin{algorithm}
\caption{Max-product belief propagation}
\begin{algorithmic}
\State Initialize messages
\Repeat
    \For{factor $a$}
        \State update $M_{a \to i}$
    \EndFor
    \For{variable $i$}
        \State update $M_{i \to a}$
    \EndFor
\Until{convergence}
\State Output $\argmax_x p(x)$ from max-marginals
\end{algorithmic}
\end{algorithm}


\subsection{Hidden Markov Models (HMMs)}

\subsubsection*{Model}

Latent chain:
\[
z_1 \to z_2 \to \cdots \to z_n, \quad x_t | z_t
\]

Factorization:
\[
p(z_{1:n}, x_{1:n})
=
p(z_1)\prod_{t=2}^n p(z_t | z_{t-1})
\prod_{t=1}^n p(x_t | z_t)
\]


\subsubsection*{Forward Pass}

Define
\[
f_t(i) = p(z_t = i, x_{1:t})
\]

Recursion:
\[
f_{t+1}(j)
=
p(x_{t+1} | z_{t+1}=j)
\sum_i f_t(i)\, p(z_{t+1}=j | z_t=i)
\]


\subsubsection*{Backward Pass}

\[
b_t(i) = p(x_{t+1:n} | z_t = i)
\]

Recursion:
\[
b_t(i)
=
\sum_j p(z_{t+1}=j | z_t=i)\, p(x_{t+1}|z_{t+1}=j)\, b_{t+1}(j)
\]


\subsubsection*{Posterior}

\[
p(z_t = i | x_{1:n})
\propto f_t(i)\, b_t(i)
\]


\subsubsection*{Viterbi Algorithm (MAP)}

Define
\[
\delta_t(j)
=
\max_{z_{1:t-1}} p(z_{1:t}, x_{1:t})
\]

Recursion:
\[
\delta_{t+1}(j)
=
p(x_{t+1}|z_{t+1}=j)
\max_i \delta_t(i)\, p(z_{t+1}=j|z_t=i)
\]

\begin{remark}
Viterbi is max-product BP on a chain.
\end{remark}


\subsection{Loopy Belief Propagation}

\begin{itemize}
\item Same updates as BP, but on graphs with cycles
\item No guarantees of convergence or correctness
\item Often works well in practice
\end{itemize}

\begin{example}
LDPC decoding uses loopy BP.
\end{example}


\subsection*{Summary: VE vs BP}

\begin{itemize}
\item \textbf{Variable elimination:}
\begin{itemize}
\item works on any graph
\item computes one marginal at a time
\item complexity depends on treewidth
\end{itemize}

\item \textbf{Belief propagation:}
\begin{itemize}
\item exact on trees
\item computes all marginals simultaneously
\item linear-time for trees
\item approximate on loopy graphs
\end{itemize}
\end{itemize}


\section{Takeaways}

\begin{itemize}
\item Exact inference is exponential in general
\item Trees $\Rightarrow$ efficient dynamic programming
\item Message passing unifies many algorithms:
\begin{itemize}
\item forward-backward
\item Viterbi
\item BP
\end{itemize}
\item MAP inference = max-product BP
\end{itemize}

\section{Sampling with Markov Chains}

\subsection*{Problem Formulation}

We want to approximately sample from a distribution
\[
p(x) \propto e^{f(x)}, \quad x \in \Omega,
\]
within $\varepsilon$ in total variation distance.

\begin{remark}
This problem appears across:
\begin{itemize}
\item Bayesian inference
\item statistical physics
\item combinatorics / theoretical CS
\item deep generative modeling
\end{itemize}
\end{remark}

\begin{remark}
Sampling is essentially equivalent to computing expectations:
\[
\E_{x \sim p}[f(x)] = \int f(x)p(x)\,dx,
\]
which becomes intractable in high dimensions.
\end{remark}


\subsection{Monte Carlo Estimation}

Given samples $X_1,\dots,X_N \sim p$,
\[
\E[f(X)] \approx \frac{1}{N} \sum_{i=1}^N f(X_i).
\]

\begin{proposition}
\[
\Var\!\left(\frac{1}{N}\sum_{i=1}^N f(X_i)\right)
=
\frac{\Var(f(X))}{N}.
\]
\end{proposition}

\begin{proof}
Uses independence:
\[
\Var\!\left(\frac{1}{N}\sum f(X_i)\right)
=
\frac{1}{N^2}\sum \Var(f(X_i))
=
\frac{1}{N}\Var(f(X)).
\]
\end{proof}


\subsection*{What Distributions Are Easy to Sample?}

\begin{itemize}
\item Uniform distributions
\item Finite discrete distributions
\item 1D distributions via inverse CDF
\item Product distributions:
\[
p(x_1,\dots,x_n) = \prod_{i=1}^n p_i(x_i)
\]
\item Directed graphical models:
\[
p(x) = \prod_{i=1}^n p(x_i | x_{\mathrm{parents}(i)})
\]
\item Gaussian distributions
\end{itemize}

\begin{remark}
General high-dimensional distributions are not directly sampleable.
\end{remark}


\subsection{Importance Sampling}

Let $Y \sim Q$.

\begin{definition}
\[
\widehat{\mu}
=
\frac{1}{N}\sum_{i=1}^N \frac{p(Y_i)}{q(Y_i)} f(Y_i).
\]
\end{definition}

\begin{proposition}
\[
\E_Q[\widehat{\mu}] = \E_P[f(X)].
\]
\end{proposition}

\begin{proof}
\[
\E_Q\left[\frac{p(Y)}{q(Y)}f(Y)\right]
=
\int \frac{p(y)}{q(y)} f(y) q(y)\,dy
=
\int f(y)p(y)\,dy.
\]
\end{proof}


\subsection{Rejection Sampling}

\begin{algorithm}
\caption{Rejection sampling}
\begin{algorithmic}
\State Input: $q(x)$, $c$ such that $p(x) \le c q(x)$
\Repeat
\State Sample $Y \sim Q$
\State Sample $U \sim \mathrm{Uniform}(0,1)$
\Until{$U \le \frac{p(Y)}{c q(Y)}$}
\State Output $Y$
\end{algorithmic}
\end{algorithm}

\begin{theorem}
If $p(x) \le c q(x)$, then rejection sampling produces samples from $p$, and
\[
\#\text{trials} \sim \mathrm{Geom}(1/c).
\]
\end{theorem}

\begin{remark}
Fails in high dimensions due to exponential scaling of $c$.
\end{remark}


\section{Markov Chains}

\begin{definition}
A sequence $(X_t)$ is a Markov chain if
\[
X_{t+1} \perp X_1,\dots,X_{t-1} \mid X_t,
\]
i.e.,
\[
p(x_{t+1} | x_1,\dots,x_t) = p(x_{t+1} | x_t).
\]
\end{definition}

\begin{definition}
A Markov chain is time-homogeneous if
\[
p(x_{t+1} | x_t) \text{ does not depend on } t.
\]
\end{definition}


\subsection*{Transition Kernels}

\begin{definition}
A Markov kernel $K$ is a function such that:
\begin{itemize}
\item $K(x,\cdot)$ is a probability measure
\item measurable in $x$
\end{itemize}
\end{definition}

For a distribution $\mu$,
\[
(\mu K)(A) = \int K(x,A)\,\mu(dx).
\]

If densities exist:
\[
(\mu K)(y) = \int \mu(x)k(x,y)\,dx.
\]


\subsection*{Operator View}

\begin{definition}
For $f : \Omega \to \R$,
\[
Kf(x) = \E_{Y \sim K(x,\cdot)}[f(Y)]
= \int f(y)K(x,dy).
\]
\end{definition}

\begin{remark}
$K$ acts:
\begin{itemize}
\item on distributions: $\mu \mapsto \mu K$
\item on functions: $f \mapsto Kf$
\end{itemize}
\end{remark}


\subsection{Stationary Distributions}

\begin{definition}
$\pi$ is stationary if
\[
\pi K = \pi.
\]
\end{definition}

\begin{remark}
Then $\pi K^t = \pi$ for all $t$.
\end{remark}


\subsection{Conditions for Convergence}

\begin{definition}
A Markov chain is:
\begin{itemize}
\item irreducible: any state can reach any other
\item aperiodic: no deterministic cycles
\end{itemize}
\end{definition}

\begin{theorem}
If a finite Markov chain is irreducible and aperiodic:
\begin{itemize}
\item it has a unique stationary distribution $\pi$
\item $\mu_0 K^t \to \pi$
\end{itemize}
\end{theorem}

\begin{theorem}[Ergodic theorem]
\[
\frac{1}{N}\sum_{t=1}^N f(X_t) \to \E_\pi[f(X)] \quad \text{a.s.}
\]
\end{theorem}


\subsection{Reversibility}

\begin{definition}
$\pi$ satisfies detailed balance if
\[
\pi(x)K(x,y) = \pi(y)K(y,x).
\]
\end{definition}

\begin{proposition}
Detailed balance $\Rightarrow$ stationarity.
\end{proposition}

\begin{remark}
Reversibility makes it easy to design Markov chains.
\end{remark}


\section{Metropolis--Hastings}

Target density: $p(x) \propto \tilde{p}(x)$.

\begin{algorithm}
\caption{Metropolis--Hastings}
\begin{algorithmic}
\State Input: $\tilde{p}(x)$, proposal $q(x,y)$
\For{$t=0,\dots,T$}
\State Sample proposal $\widehat{X}_{t+1} \sim q(X_t,\cdot)$
\State Accept with probability
\[
\alpha(x,y)
=
\min\left\{
\frac{\tilde{p}(y)q(y,x)}{\tilde{p}(x)q(x,y)}, 1
\right\}
\]
\EndFor
\end{algorithmic}
\end{algorithm}

\begin{theorem}
The MH chain has stationary distribution $p$.
\end{theorem}


\subsection*{Gibbs Sampling}

\begin{definition}
At each step:
\[
X_i \sim p(x_i | x_{-i}).
\]
\end{definition}

\begin{proposition}
Gibbs sampling is a special case of Metropolis--Hastings with acceptance probability $1$.
\end{proposition}

\begin{remark}
In graphical models:
\[
p(x_i | x_{-i}) = p(x_i | x_{\mathrm{MB}(i)}).
\]
\end{remark}


\subsection*{Examples}

\begin{example}[Ball walk]
\[
\widehat{X}_{t+1} \sim \mathrm{Uniform}(B_R(X_t)),
\]
then accept/reject via MH.
\end{example}

\begin{example}[Gibbs sampling]
Select coordinate $i$ and sample
\[
X_i \sim p(x_i | x_{-i}).
\]
\end{example}


\section*{Takeaways}

\begin{itemize}
\item Direct sampling is rarely possible
\item MCMC constructs a Markov chain with target stationary distribution
\item Convergence requires irreducibility + aperiodicity
\item Reversibility simplifies construction (MH)
\item Gibbs exploits conditional structure
\end{itemize}

\section{Markov Chain Monte Carlo (MCMC)}

\subsection{Sampling via Markov Chains}

\begin{definition}
Goal: sample from a distribution
\[
p(x) \propto e^{f(x)}, \quad x \in \Omega,
\]
within $\varepsilon$ in total variation distance.
\end{definition}

\begin{remark}
We construct a Markov chain whose stationary distribution is $p$, then simulate it.
\end{remark}


\subsection*{Rejection Sampling (Limitations)}

\begin{theorem}
If $p(x) \le c q(x)$, rejection sampling produces samples from $p$, and
\[
\#\text{trials} \sim \mathrm{Geom}(1/c).
\]
\end{theorem}

\begin{remark}
In high dimensions, $c$ typically grows exponentially, making rejection sampling impractical.
\end{remark}


\subsection*{Markov Chains}

\begin{definition}
A sequence $(X_t)$ is a Markov chain if
\[
p(x_{t+1} | x_1,\dots,x_t) = p(x_{t+1} | x_t).
\]
\end{definition}

\begin{definition}
A distribution $\pi$ is stationary if
\[
\pi K = \pi.
\]
\end{definition}


\subsection{Reversibility}

\begin{definition}
$\pi$ satisfies detailed balance if
\[
\pi(x)K(x,y) = \pi(y)K(y,x).
\]
\end{definition}

\begin{proposition}
If detailed balance holds, then $\pi$ is stationary.
\end{proposition}

\begin{proof}
\[
\sum_x \pi(x)K(x,y)
=
\sum_x \pi(y)K(y,x)
=
\pi(y)\sum_x K(y,x)
=
\pi(y).
\]
\end{proof}


\section*{Metropolis--Hastings}

\begin{algorithm}
\caption{Metropolis--Hastings}
\begin{algorithmic}
\State Input: $\tilde{p}(x)$, proposal $q(x,y)$
\For{$t=0,\dots,T$}
\State Sample $Y \sim q(X_t,\cdot)$
\State Accept with probability
\[
\alpha(x,y)
=
\min\left\{
\frac{\tilde{p}(y)q(y,x)}{\tilde{p}(x)q(x,y)}, 1
\right\}
\]
\State Set $X_{t+1} = Y$ if accepted, else $X_{t+1} = X_t$
\EndFor
\end{algorithmic}
\end{algorithm}

\begin{theorem}
The Metropolis--Hastings chain has stationary distribution $p$.
\end{theorem}

\begin{proof}
Check detailed balance:
\[
p(x)K(x,y)
=
p(x)q(x,y)\alpha(x,y)
=
\min\{p(x)q(x,y), p(y)q(y,x)\},
\]
which is symmetric in $(x,y)$.
\end{proof}


\subsection*{Gibbs Sampling}

\begin{definition}
At each step, pick coordinate $i$ and sample
\[
X_i \sim p(x_i | x_{-i}).
\]
\end{definition}

\begin{proposition}
Gibbs sampling is a special case of Metropolis--Hastings with acceptance probability $1$.
\end{proposition}

\begin{remark}
In graphical models:
\[
p(x_i | x_{-i}) = p(x_i | x_{\mathrm{MB}(i)}).
\]
\end{remark}


\section{Langevin Monte Carlo (LMC)}

\subsection*{Continuous-time process}

\begin{definition}
Langevin diffusion:
\[
dX_t = -\nabla V(X_t)\,dt + \sqrt{2}\,dB_t.
\]
\end{definition}

\begin{remark}
Target distribution:
\[
p(x) \propto e^{-V(x)}.
\]
\end{remark}


\subsubsection*{Discretization}

Using Euler--Maruyama:
\[
X_{t+1} = X_t - \eta \nabla V(X_t) + \sqrt{2\eta}\,\xi_t,
\quad \xi_t \sim \Normal(0,I).
\]

\begin{remark}
This is noisy gradient descent.
\end{remark}


\subsection*{Stationary Distribution}

\begin{theorem}
The Langevin diffusion has stationary distribution
\[
p(x) \propto e^{-V(x)}.
\]
\end{theorem}

\begin{proof}
From the Fokker--Planck equation:
\[
\frac{\partial \mu_t}{\partial t}
=
\nabla \cdot (\nabla V \,\mu_t) + \Delta \mu_t.
\]

Plug $\mu(x) = e^{-V(x)}$:

\[
\nabla \mu = -(\nabla V)e^{-V}, \quad
\Delta \mu = \left(-\Delta V + \|\nabla V\|^2\right)e^{-V}.
\]

Then:
\[
\nabla \cdot (\nabla V \mu)
=
(\Delta V)e^{-V} - \|\nabla V\|^2 e^{-V}.
\]

Adding:
\[
\nabla \cdot (\nabla V \mu) + \Delta \mu = 0.
\]

Thus $\mu$ is stationary.
\end{proof}


\subsection{MALA (Metropolis-adjusted Langevin)}

\begin{algorithm}
\caption{MALA}
\begin{algorithmic}
\State Propose
\[
Y \sim \Normal(x - h \nabla V(x), 2hI)
\]
\State Accept with MH correction
\end{algorithmic}
\end{algorithm}

\begin{remark}
Corrects discretization bias of LMC.
\end{remark}


\section{Underdamped Langevin Dynamics}

\begin{definition}
Introduce momentum $v$:
\[
\begin{aligned}
dX_t &= v_t\,dt \\
dv_t &= -\gamma v_t\,dt - \nabla V(X_t)\,dt + \sqrt{2\gamma}\,dB_t
\end{aligned}
\]
\end{definition}

\begin{remark}
Stationary distribution:
\[
p(x,v) \propto e^{-V(x) - \frac{1}{2}\|v\|^2}.
\]
\end{remark}

\begin{remark}
Non-reversible → faster mixing.
\end{remark}


\section*{Hamiltonian Monte Carlo (HMC)}

\subsection*{Hamiltonian dynamics}

Define:
\[
H(q,p) = U(q) + K(p)
\]

\[
\begin{aligned}
\frac{dq}{dt} &= \nabla_p H = \nabla K(p) \\
\frac{dp}{dt} &= -\nabla_q H = -\nabla U(q)
\end{aligned}
\]

\begin{proposition}
Hamiltonian dynamics preserve:
\begin{itemize}
\item Energy: $H(q(t),p(t))$
\item Volume (Liouville's theorem)
\end{itemize}
\end{proposition}


\subsection*{Reversibility trick}

Define momentum flip:
\[
R(q,p) = (q,-p).
\]

Then:
\[
R \circ S_T \circ R = S_T^{-1}.
\]

\begin{proposition}
HMC transition is reversible w.r.t.
\[
\mu(q,p) \propto e^{-H(q,p)}.
\]
\end{proposition}


\subsection*{Algorithm}

\begin{algorithm}
\caption{Hamiltonian Monte Carlo (idealized)}
\begin{algorithmic}
\State Sample $p \sim e^{-K(p)}$
\State Simulate Hamiltonian dynamics for time $T$
\State Output $q$
\end{algorithmic}
\end{algorithm}

\begin{remark}
No rejection needed in exact dynamics (but required after discretization).
\end{remark}


\subsection*{Leapfrog Integrator}

\begin{algorithm}
\caption{Leapfrog}
\begin{algorithmic}
\State $p \leftarrow p - \frac{h}{2}\nabla U(q)$
\For{$t=1,\dots,T$}
\State $q \leftarrow q + h\nabla K(p)$
\State $p \leftarrow p - h\nabla U(q)$
\EndFor
\State $p \leftarrow p - \frac{h}{2}\nabla U(q)$
\end{algorithmic}
\end{algorithm}

\begin{remark}
Properties:
\begin{itemize}
\item volume-preserving
\item reversible
\item approximately energy-preserving
\end{itemize}
\end{remark}


\section*{Takeaways}

\begin{itemize}
\item MCMC constructs Markov chains with target stationary distribution
\item Metropolis--Hastings ensures correctness via detailed balance
\item Gibbs exploits conditional structure
\item Langevin adds gradient information (diffusion-based sampling)
\item HMC uses Hamiltonian dynamics for efficient exploration
\end{itemize}

\subsection*{Langevin Monte Carlo}

We want to sample from
\[
p(x) \propto e^{-V(x)}.
\]

\begin{definition}
The Langevin diffusion is the stochastic differential equation
\[
dX_t = -\nabla V(X_t)\,dt + \sqrt{2}\,dB_t.
\]
\end{definition}

Discretizing via Euler--Maruyama:
\[
X_{t+1} = X_t - \eta \nabla V(X_t) + \sqrt{2\eta}\,\xi_t,
\quad \xi_t \sim \Normal(0,I_d).
\]

\begin{remark}
This can be viewed as gradient descent with injected Gaussian noise.
\end{remark}

\begin{remark}
The continuous-time process has stationary distribution $p(x) \propto e^{-V(x)}$, but the discretized version introduces bias unless $\eta$ is small.
\end{remark}


\subsection*{Metropolis-Adjusted Langevin Algorithm (MALA)}

To correct discretization bias, use Langevin as a proposal in Metropolis--Hastings.

\begin{definition}
Proposal distribution:
\[
q(x,y)
\propto
\exp\left(
-\frac{\|y - (x - \eta \nabla V(x))\|^2}{4\eta}
\right).
\]
\end{definition}

\begin{algorithm}
\caption{MALA}
\begin{algorithmic}
\State Propose $Y \sim \Normal(x - \eta \nabla V(x),\, 2\eta I)$
\State Accept with probability
\[
\alpha(x,y)
=
\min\left\{
\frac{p(y)q(y,x)}{p(x)q(x,y)}, 1
\right\}
\]
\end{algorithmic}
\end{algorithm}


\subsection*{Underdamped Langevin Dynamics}

\begin{definition}
Introduce momentum $v$:
\[
\begin{aligned}
dX_t &= v_t\,dt, \\
dv_t &= -\gamma v_t\,dt - \nabla V(X_t)\,dt + \sqrt{2\gamma}\,dB_t.
\end{aligned}
\]
\end{definition}

\begin{remark}
Stationary distribution:
\[
p(x,v) \propto e^{-V(x) - \frac{1}{2}\|v\|^2}.
\]
\end{remark}

\begin{remark}
Noise is injected only in $v_t$, leading to more efficient exploration.
\end{remark}


\subsection*{Continuous-Time Markov Processes}

\begin{definition}
A continuous-time Markov process $(X_t)_{t \ge 0}$ satisfies
\[
P(X_u \in A | X_{[0,t]}) = P(X_u \in A | X_t), \quad u > t.
\]
\end{definition}

\begin{definition}
A family $(P_t)_{t \ge 0}$ is a Markov semigroup if
\[
P_{t+s} = P_t P_s, \quad P_0 = \mathrm{id}.
\]
\end{definition}


\subsection*{Generator of a Markov Process}

\begin{definition}
The generator $\mathcal{L}$ is
\[
\mathcal{L}f = \left|\frac{d}{dt} P_t f \right|_{t=0}.
\]
\end{definition}

\begin{remark}
This defines the process via
\[
P_t = e^{t\mathcal{L}}.
\]
\end{remark}


\subsection{Fokker--Planck Equation}

Let $\mathcal{L}^*$ be the adjoint operator.

\begin{definition}
The distribution $\mu_t$ evolves as
\[
\frac{d\mu_t}{dt} = \mathcal{L}^* \mu_t.
\]
\end{definition}

\begin{remark}
For Langevin diffusion:
\[
\mathcal{L}f = -\langle \nabla f, \nabla V\rangle + \Delta f,
\]
\[
\mathcal{L}^* \mu = \nabla \cdot (\nabla V \mu) + \Delta \mu.
\]
\end{remark}


\subsection*{Stationarity of Langevin}

\begin{theorem}
The distribution
\[
\mu(x) \propto e^{-V(x)}
\]
is stationary for Langevin diffusion.
\end{theorem}

\begin{proof}
Plug into Fokker--Planck equation:
\[
\mathcal{L}^*\mu
=
\nabla \cdot (\nabla V \mu) + \Delta \mu.
\]

Compute:
\[
\nabla \mu = -(\nabla V)\mu,
\quad
\Delta \mu = (-\Delta V + \|\nabla V\|^2)\mu.
\]

Then:
\[
\nabla \cdot (\nabla V \mu)
=
(\Delta V)\mu - \|\nabla V\|^2 \mu.
\]

Adding:
\[
(\Delta V - \|\nabla V\|^2)\mu + (-\Delta V + \|\nabla V\|^2)\mu = 0.
\]
\end{proof}


\section{Hamiltonian Monte Carlo (HMC)}

\subsection{Hamiltonian Dynamics}

Define
\[
H(q,p) = U(q) + K(p).
\]

Dynamics:
\[
\frac{dq}{dt} = \nabla_p H, \quad
\frac{dp}{dt} = -\nabla_q H.
\]

\begin{remark}
Typical choice:
\[
K(p) = \frac{1}{2}\|p\|^2.
\]
\end{remark}


\subsection{Key Properties}

\begin{proposition}
Hamiltonian dynamics preserve:
\begin{itemize}
\item energy $H(q,p)$
\item volume in phase space
\end{itemize}
\end{proposition}


\subsection{Leapfrog Integrator}

\begin{lemma}
The leapfrog integrator:
\begin{itemize}
\item preserves volume
\item is reversible (up to momentum flip)
\end{itemize}
\end{lemma}

\begin{remark}
It does \emph{not} exactly preserve the Hamiltonian.
\end{remark}


\subsection*{HMC with Leapfrog}

\begin{algorithm}
\caption{Hamiltonian Monte Carlo (practical)}
\begin{algorithmic}
\State Sample $p \sim e^{-K(p)}$
\State Simulate Hamiltonian dynamics via leapfrog for time $T$
\State Accept with probability
\[
\alpha = \min\{e^{-H(q',p') + H(q,p)}, 1\}
\]
\end{algorithmic}
\end{algorithm}

\begin{remark}
The Metropolis step corrects discretization error.
\end{remark}


\subsection*{Choosing Parameters}

\begin{itemize}
\item Integration time $T$:
\begin{itemize}
\item Too small $\Rightarrow$ slow exploration
\item Too large $\Rightarrow$ wasted computation
\end{itemize}
\item Randomizing $T$ helps avoid periodic behavior
\item No-U-Turn criterion adaptively selects trajectory length
\end{itemize}


\subsection*{Geometric Adaptation}

\begin{definition}
Riemannian HMC uses position-dependent metric:
\[
K(q,p) = \frac{1}{2} p^\top \Sigma^{-1}(q)p + \frac{1}{2}\log \det \Sigma(q).
\]
\end{definition}


\section{Convergence of Markov Chains}

\subsection{Spectral Gap}

\begin{theorem}
Let $\lambda_2$ be the second largest eigenvalue. Then
\[
\Var_\mu(K^n f) \le |\lambda_2|^{2n} \Var_\mu(f).
\]
\end{theorem}


\subsection{Non-reversible Chains}

\begin{remark}
For non-reversible chains:
\[
\chi^2(\nu K^n || \mu)
\le C |\lambda_2|^{2n} \chi^2(\nu || \mu).
\]
\end{remark}

\begin{remark}
In continuous time this is called \emph{hypocoercivity}.
\end{remark}


\subsubsection{Poincar\'e Inequality}

\begin{definition}
A distribution $\mu$ satisfies a Poincar\'e inequality if
\[
\langle g, \mathcal{L}g \rangle_\mu \ge c \Var_\mu(g).
\]
\end{definition}

\begin{remark}
Implies exponential convergence rate $e^{-2ct}$.
\end{remark}


\subsection{Contraction in Divergences}

\begin{definition}
A kernel $K$ is contractive if
\[
D_f(\nu K || \mu) \le (1-\kappa) D_f(\nu || \mu).
\]
\end{definition}

\begin{remark}
Implies geometric convergence.
\end{remark}


\subsection{Mixing Time}

\begin{definition}
\[
t_{\mathrm{mix}}(\varepsilon)
=
\min\{t : \sup_{\mu_0} \mathrm{TV}(\mu_0 K^t, \mu) \le \varepsilon\}.
\]
\end{definition}

\begin{definition}
Relaxation time:
\[
t_{\mathrm{rel}} = \frac{1}{1-|\lambda_2|}.
\]
\end{definition}


\subsection{Conductance}

\begin{definition}
For $S \subset \Omega$:
\[
\Phi(S) = \frac{Q(S,S^c)}{\mu(S)},
\quad
Q(S,S^c) = \sum_{x \in S, y \in S^c} \mu(x)P(x,y).
\]
\end{definition}

\begin{definition}
\[
\Phi_* = \min_{\mu(S)\le 1/2} \Phi(S).
\]
\end{definition}

\begin{theorem}
\[
t_{\mathrm{mix}}(1/4) \ge \frac{1}{4\Phi_*}.
\]
\end{theorem}


\subsection{Asymptotic Variance}

\begin{definition}
\[
\sigma_f^2
=
\Var_\mu(f(X_0)) + 2\sum_{k=1}^\infty \mathrm{Cov}_\mu(f(X_0), f(X_k)).
\]
\end{definition}

\begin{theorem}
\[
\frac{1}{N}\sum_{k=0}^{N-1} f(X_k)
\to \E_\mu[f]
\]
and
\[
\Var\left(\frac{1}{N}\sum f(X_k)\right) \sim \frac{\sigma_f^2}{N}.
\]
\end{theorem}


\subsection{Practical Diagnostics}

\begin{itemize}
\item Trace plots (check stationarity)
\item Autocorrelation:
\[
\rho_\ell = \frac{\mathrm{Cov}(f(X_0), f(X_\ell))}{\Var(f)}
\]
\item Effective sample size:
\[
N_{\mathrm{eff}} = \frac{N}{1 + 2\sum_{\ell \ge 1} \rho_\ell}
\]
\end{itemize}


\subsection*{Discussion}

\begin{itemize}
\item Single long chain vs multiple shorter chains:
\begin{itemize}
\item single chain: cheaper burn-in
\item multiple chains: better exploration of multimodal distributions
\end{itemize}
\end{itemize}

\section{Tempering and Particle Methods}

\subsection{Multimodality and Slow Mixing}

Multimodal target distributions pose a fundamental challenge for MCMC.

\begin{itemize}
    \item Probability mass is concentrated in separated regions (modes).
    \item Transitions between modes require crossing \emph{energy barriers}.
    \item Local proposals (e.g., random walk, Langevin) rarely cross these barriers.
\end{itemize}

\begin{remark}
For distributions of the form $p(x) \propto e^{-f(x)}$, the time to transition between modes is often \emph{exponential} in the barrier height:
\[
\text{mixing time} \sim \exp(\Delta f).
\]
This is known as \emph{metastability}.
\end{remark}

\subsection{Temperature Scaling}

Introduce a \emph{temperature parameter} $T > 0$:
\[
p_T(x) \propto e^{-f(x)/T} = e^{-\beta f(x)}, \quad \beta = \frac{1}{T}.
\]

\begin{itemize}
    \item High temperature ($T \gg 1$, $\beta \ll 1$): distribution is \emph{flattened}.
    \item Low temperature ($T \ll 1$, $\beta \gg 1$): distribution concentrates near modes.
\end{itemize}

\begin{remark}
In Langevin dynamics:
\[
dX_t = -\nabla V(X_t)\,dt + \sqrt{2T}\,dB_t,
\]
temperature controls the tradeoff between gradient (drift) and noise.
\end{remark}

\subsection{Simulated Annealing}

Simulated annealing gradually lowers temperature:
\[
T_1 > T_2 > \cdots \to 0.
\]

\begin{itemize}
    \item Initially explores broadly.
    \item Eventually concentrates near global minima.
\end{itemize}

\begin{remark}
Once modes separate at low temperature, the chain becomes \emph{trapped}:
probability mass cannot easily move between modes.
\end{remark}

\subsection{Simulated Tempering}

Introduce temperature as an auxiliary variable:
\[
(\ell, x) \in \{1,\dots,L\} \times \Omega.
\]

Define target:
\[
\pi(\ell, x) \propto w_\ell\, p_\ell(x), \quad p_\ell(x) \propto e^{-f(x)/T_\ell}.
\]

Algorithm:
\begin{itemize}
    \item Alternate:
    \begin{itemize}
        \item Update $x$ using MCMC at fixed temperature $\ell$.
        \item Propose $\ell \to \ell \pm 1$ with MH acceptance:
        \[
        \alpha = \min\left\{ \frac{w_{\ell'} p_{\ell'}(x)}{w_\ell p_\ell(x)}, 1 \right\}.
        \]
    \end{itemize}
\end{itemize}

\begin{remark}
Drawback: requires knowledge (or estimation) of normalizing constants for $p_\ell$.
\end{remark}

\subsection{Parallel Tempering (Replica Exchange)}

Run multiple chains at different temperatures:
\[
(x_1, \dots, x_L), \quad x_\ell \sim p_\ell.
\]

Target:
\[
\pi(x_1,\dots,x_L) = \prod_{\ell=1}^L p_\ell(x_\ell).
\]

Swap proposal:
\[
(x_i, x_{i+1}) \to (x_{i+1}, x_i),
\]
with acceptance probability
\[
\alpha = \min\left\{ \frac{p_i(x_{i+1}) p_{i+1}(x_i)}{p_i(x_i)p_{i+1}(x_{i+1})}, 1 \right\}.
\]

\begin{remark}
Advantages:
\begin{itemize}
    \item No need for normalizing constants.
    \item High-temperature chains help low-temperature chains escape modes.
\end{itemize}
\end{remark}

\subsection{Sequential Monte Carlo (SMC)}

Construct a sequence of distributions:
\[
p_k(x) \propto \tilde p_k(x), \quad k = 1,\dots,L,
\]
typically interpolating between easy and target distributions.

\begin{algorithm}[H]
\caption{Sequential Monte Carlo}
\begin{algorithmic}[1]
\State Sample $X_1^{(i)} \sim p_1$
\For{$k = 2$ to $L$}
    \State Compute weights:
    \[
    w_k^{(i)} \propto \frac{\tilde p_k(X_{k-1}^{(i)})}{\tilde p_{k-1}(X_{k-1}^{(i)})}
    \]
    \State Resample particles according to $w_k^{(i)}$
    \State Move particles: $X_k^{(i)} \sim K_k(X_{k-1}^{(i)}, \cdot)$
\EndFor
\State Output $\{X_L^{(i)}\}$
\end{algorithmic}
\end{algorithm}

\begin{remark}
SMC combines:
\begin{itemize}
    \item Importance sampling (weights)
    \item Resampling (to prevent degeneracy)
    \item MCMC moves (to diversify particles)
\end{itemize}
\end{remark}

\subsection{Annealed Importance Sampling (AIS)}

A simplified version of SMC without resampling.

\begin{algorithm}[H]
\caption{Annealed Importance Sampling}
\begin{algorithmic}[1]
\State Sample $X_1^{(i)} \sim p_1$
\For{$k = 2$ to $L$}
    \State Sample $X_k^{(i)} \sim K_k(X_{k-1}^{(i)}, \cdot)$
\EndFor
\State Compute weights:
\[
\tilde w^{(i)} = \prod_{k=2}^L \frac{\tilde p_k(X_k^{(i)})}{\tilde p_{k-1}(X_k^{(i)})}
\]
\State Estimate:
\[
\E_{p_L}[f(X)] \approx \frac{\sum_i \tilde w^{(i)} f(X_L^{(i)})}{\sum_i \tilde w^{(i)}}
\]
\end{algorithmic}
\end{algorithm}

\begin{remark}
AIS is especially useful for:
\begin{itemize}
    \item Estimating partition functions
    \item Bridging between simple and complex distributions
\end{itemize}
\end{remark}

\subsection*{Key Insight}

Tempering and particle methods address the same core issue:

\begin{center}
\emph{Exploration vs. concentration tradeoff}
\end{center}

\begin{itemize}
    \item High temperature $\Rightarrow$ exploration
    \item Low temperature $\Rightarrow$ accurate sampling
\end{itemize}

\begin{itemize}
    \item Tempering methods: move between temperatures
    \item Particle methods: propagate distributions across temperatures
\end{itemize}



\end{document}